\begin{lemma}{Kompaktheit der Sph�re}{PraekompaktheitDerSphaere}
Sei $(E, \norm\cdot_E)$ ein normierter $\K$--Vektorraum.\\
Dann gilt:
\begin{enumerate}[(1)]
	\item Ist $\kran E01$ kompakt, so auch $\kabg E01$.
	\item Ist $(F, \norm\cdot_F)$ ein weiterer normierter $\K$--Vektorraum,
	      $\kran {L(E,F)}01$ kompakt und $\Inv(E,F) \neq \emptyset$, so sind
	      $E$ und $F$ endlichdimensional.
\end{enumerate}
\end{lemma}
\begin{proof}
(1) Die Aussage ist trivial, wenn $E = \meng 0$.\\
Sei also zus�tzlich $E \neq \meng 0$ angenommen.\\
Dann ist $\kran E 01 \neq \emptyset$ und offensichtlich gilt
\begin{align}\kabg E 01 = \menge{sy}{s \in [0,1], y \in \kran E 01}\text. 
\tag {$\ast$}\label{pdseq:1}\end{align}
Sei nun $K:=[0,1] \times \kran E01$ mit der Produktmetrik ausgestattet.\\
Ist $\kran E01$ kompakt, so ist auch $K$ kompakt.\\
Seien $f: K \rightarrow \kabg E 01, (s,y) \mapsto sy$. Dann ist $f$ stetig.\\
Wegen $\eqref{pdseq:1}$ ist $f(K)=\kabg E01$ und damit ist $\kabg E01$ als
stetiges Bild einer kompakten Menge selbst kompakt.\\
(3) Sei $(F, \norm\cdot_F)$ ein normierter $\K$--Vektorraum und es gelte
$\kran {L(E,F)}01$ ist kompakt und $\Inv(E, F) \neq \emptyset$.\\
Nach (2) ist dann $\kabg{L(E,F)}01$ kompakt, also $L(E,F)$
endlichdimensional.\\
Ist $F= \meng 0$, so ist wegen $\Inv(E,F) \neq \emptyset$ auch $E = \meng 0$.\\
Ist $F\neq0$, so existiert nach Hahn--Banach ein $\vi\in F'\setminus\meng 0$.\\
Sei $\Phi: F \rightarrow L(E,F), y \mapsto (x \mapsto \vi(x)y)$.\\
Offensichtlich ist $\Phi$ linear und wegen $\vi \neq 0$ auch injektiv.\\
Damit ist $F$ isomorph zu einem Unterraum des endlichdimensionalen Raums
$L(E,F)$, also endlichdimensional und
wegen $\Inv(E,F) \neq \emptyset$, ist es auch $E$.
\end{proof}